\documentclass[a4paper,12pt]{article}

\usepackage[italian]{babel}
\usepackage[utf8]{inputenc}
\usepackage{geometry}
\usepackage{amsmath}
\usepackage{array}
\usepackage{circuitikz}

\geometry{margin=2.5cm}

\title{Verifica di Informatica}
\author{Prof. Fedeli Massimo}
\date{}

\begin{document}
	
	\maketitle
	
	Nome e Cognome: \rule{8cm}{0.4pt}
	
	Classe: \rule{3cm}{0.4pt} \hfill Data: \rule{3cm}{0.4pt}
	
	\vspace{1cm}
	
	\section*{Esercizio 1 – Algoritmi e Flowgorithm}
	
	Scrivi un algoritmo in \textbf{Flowgorithm} che:
	
	\begin{itemize}
		\item chieda all'utente quanti numeri interi vuole inserire;
		\item il numero di valori deve essere \textbf{almeno 10}. Se l'utente inserisce un numero minore a 10, l'algoritmo deve richiedere nuovamente il valore;
		\item memorizzi i valori inseriti in un \textbf{vettore di interi};
		\item al termine dell'inserimento determini:
		\begin{itemize}
			\item la \textbf{somma} dei valori
			\item la \textbf{media}
			\item il \textbf{valore massimo}
			\item il \textbf{valore minimo}
			\item il \textbf{numero che compare più volte} nel vettore
		\end{itemize}
	\end{itemize}
	
	L'algoritmo deve essere corredato da \textbf{commenti esplicativi} che descrivano le principali operazioni svolte.
	
	\vspace{1cm}
	
	\section*{Esercizio 2 – Funzioni logiche}
	
	Data la seguente \textbf{tabella di verità} con tre variabili logiche \(A\), \(B\) e \(C\):
	
	\begin{center}
		\begin{tabular}{|c|c|c|c|}
			\hline
			A & B & C & F \\ 
			\hline
			0 & 0 & 0 & 0 \\ 
			0 & 0 & 1 & 1 \\ 
			0 & 1 & 0 & 0 \\ 
			0 & 1 & 1 & 1 \\ 
			1 & 0 & 0 & 1 \\ 
			1 & 0 & 1 & 1 \\ 
			1 & 1 & 0 & 0 \\ 
			1 & 1 & 1 & 1 \\ 
			\hline
		\end{tabular}
	\end{center}
	
	Determinare:
	
	\begin{enumerate}
		\item la \textbf{funzione logica} \(F(A,B,C)\) in forma canonica (somma di prodotti);
		\item la \textbf{semplificazione} della funzione;
		\item il \textbf{circuito logico equivalente} utilizzando porte AND, OR e NOT.
	\end{enumerate}
	
	\vspace{1cm}
	
	\section*{Esercizio 3 – Analisi di circuito logico}
	
	Dato il seguente circuito logico:
	
	\begin{center}
		\begin{circuitikz}
			
			\node (A) at (0,2) {$A$};
			\node (B) at (0,0) {$B$};
			\node (C) at (0,-2) {$C$};
			
			\draw (A) -- (1,2);
			\draw (B) -- (1,0);
			\draw (C) -- (1,-2);
			
			\draw (1,2) -- (2,2) node[and port, anchor=in 1] (and1) {};
			\draw (1,0) -- (and1.in 2);
			
			\draw (1,-2) -- (2,-2) node[not port] (not1) {};
			
			\draw (and1.out) -- (4,1);
			\draw (not1.out) -- (4,-1);
			
			\node[or port] (or1) at (5,0) {};
			
			\draw (4,1) -- (or1.in 1);
			\draw (4,-1) -- (or1.in 2);
			
			\draw (or1.out) -- (6,0) node[right] {$F$};
			
		\end{circuitikz}
	\end{center}
	
	Determinare:
	
	\begin{enumerate}
		\item la \textbf{funzione logica} realizzata dal circuito;
		\item la \textbf{tabella di verità} completa;
		\item indicare se la funzione può essere ulteriormente \textbf{semplificata}.
	\end{enumerate}
	
	\vspace{1cm}
	
	\textbf{Tempo a disposizione:} 60 minuti
	
	\textbf{È consentito l'uso di fogli di brutta copia.}
	
\end{document}